% =====================================================================
% Chapter: Maxwell's equations, the wave equation, and the boundary
%          terms we usually discard
%
% Requires: amsmath, amssymb; a class with \chapter (book/report).
%
% Conventions
%   Sections 1-3 : SI units.
%   Sections 4-8 : Heaviside-Lorentz units, c = eps_0 = mu_0 = 1,
%                  metric (+,-,-,-), A^mu = (phi, A),
%                  F^{0i} = -E^i,  F^{ij} = -eps^{ijk} B^k.
% =====================================================================

\chapter[Maxwell's equations and boundary terms]{Maxwell's equations, the wave equation, and the boundary terms we usually discard}
\label{ch:maxwell-boundary}

%----------------------------------------------------------------------
\section[Two textbook statements]{Two textbook statements that do not sit comfortably together}
\label{sec:mb-intro}

Every textbook on electrodynamics contains the following two statements.

\begin{enumerate}
\item Maxwell's equations imply that the fields (and, in a suitable gauge,
the potentials) satisfy the wave equation. Without sources this equation
still has non-trivial solutions: light.
\item Maxwell's equations follow from a least-action principle. One varies
the potentials, integrates by parts, ``neglects the total divergence'',
and reads off the field equations.
\end{enumerate}

Both statements are correct. Placed side by side, however, they hide a
tension. A source-free solution of a linear hyperbolic equation can be
different from zero only if \emph{something} is different from zero on the
boundary of the region in which we solve it. The variational principle, on
the other hand, is derived by throwing the boundary away.

The thesis of this chapter is that these are \emph{one fact seen from two
sides}. The total divergences that we discard in the action are precisely
the place where the source-free solutions of the wave equation---what
quantum theory calls photons---are visible to the action at all. Discarding
them is harmless for obtaining the field equations. It is not harmless for
what the field equations \emph{cannot} tell us: which free radiation is
present, what the symplectic structure of the theory is, and in which
direction time runs.

The plan is as follows. Sections~\ref{sec:mb-maxwell} and~\ref{sec:mb-wave}
recall Maxwell's equations and derive the wave equation, paying attention
to what is lost on the way. Section~\ref{sec:mb-anatomy} splits every
solution into a particular and a homogeneous part and shows that the
homogeneous part is controlled by boundary data. Section~\ref{sec:mb-action}
turns to the action and examines exactly which divergences are dropped and
why. Sections~\ref{sec:mb-symplectic} and~\ref{sec:mb-arrow} show what
those divergences contain: the symplectic structure (hence the photon
number) and the arrow of time (hence radiation reaction).

%----------------------------------------------------------------------
\section{Maxwell's equations}
\label{sec:mb-maxwell}

In SI units, with $c^2=1/(\varepsilon_0\mu_0)$,
\begin{align}
\nabla\cdot\mathbf E &= \rho/\varepsilon_0, &
\nabla\times\mathbf B-\frac1{c^2}\,\partial_t\mathbf E &= \mu_0\mathbf J,
\label{eq:mb-max-a}\\
\nabla\cdot\mathbf B &= 0, &
\nabla\times\mathbf E+\partial_t\mathbf B &= 0 .
\label{eq:mb-max-b}
\end{align}
There are eight scalar equations for six unknown field components, so the
system is overdetermined. It is nevertheless consistent, for two reasons.

First, the divergence equations are conditions on the \emph{data}, not on
the dynamics. Taking the divergence of the two curl equations and using
the divergence equations gives
\begin{equation}
\partial_t\rho+\nabla\cdot\mathbf J=0,\qquad
\partial_t(\nabla\cdot\mathbf B)=0,
\label{eq:mb-consistency}
\end{equation}
so if the constraints $\nabla\cdot\mathbf E=\rho/\varepsilon_0$ and
$\nabla\cdot\mathbf B=0$ hold at one instant, and charge is conserved, they
hold at all later instants. Second, only the two curl equations contain time
derivatives; they are the evolution equations.

Charge conservation is not an additional assumption: it is a necessary
condition for the sourced equations to have any solution at all.

%----------------------------------------------------------------------
\section{From Maxwell's equations to the wave equation}
\label{sec:mb-wave}

\subsection{Wave equations for the fields}

Apply $\nabla\times$ to Faraday's law and use
$\nabla\times(\nabla\times\mathbf E)=\nabla(\nabla\cdot\mathbf E)-\nabla^2\mathbf E$
together with Amp\`ere's law:
\begin{equation}
\Box\,\mathbf E=-\frac1{\varepsilon_0}\nabla\rho-\mu_0\,\partial_t\mathbf J,
\qquad
\Box\,\mathbf B=\mu_0\,\nabla\times\mathbf J,
\qquad
\Box\equiv\frac1{c^2}\partial_t^2-\nabla^2 .
\label{eq:mb-wave-fields}
\end{equation}
The second equation follows in the same way from Amp\`ere's law and
$\nabla\cdot\mathbf B=0$.

These equations are a \emph{consequence} of Maxwell's equations, not an
equivalent reformulation. Differentiating raises the order, and the extra
solutions are spurious. For example, $\mathbf E=\hat{\mathbf x}\cos(kz-\omega t)$
and $\mathbf B=\hat{\mathbf x}\cos(kz-\omega t)/c$ solve the two homogeneous
wave equations separately, but they violate Faraday's law, because
$\nabla\times\mathbf E$ points along $\hat{\mathbf y}$ while $\partial_t\mathbf B$
points along $\hat{\mathbf x}$. Two copies of a wave equation know nothing
about each other; it is the first-order system that couples $\mathbf E$ to
$\mathbf B$. Whenever we use the wave equation we must remember to impose
the first-order equations on the data.

\subsection{Potentials and gauge}

Equations~\eqref{eq:mb-max-b} are solved identically by
\begin{equation}
\mathbf B=\nabla\times\mathbf A,\qquad
\mathbf E=-\nabla\phi-\partial_t\mathbf A .
\label{eq:mb-potentials}
\end{equation}
The remaining equations~\eqref{eq:mb-max-a} become
\begin{align}
-\nabla^2\phi-\partial_t(\nabla\cdot\mathbf A)&=\rho/\varepsilon_0,\\
\Box\mathbf A+\nabla\!\left(\nabla\cdot\mathbf A+\frac1{c^2}\partial_t\phi\right)&=\mu_0\mathbf J .
\end{align}
The potentials are not unique: $\mathbf A\to\mathbf A+\nabla\chi$,
$\phi\to\phi-\partial_t\chi$ leaves $\mathbf E$ and $\mathbf B$ unchanged. In
the Lorenz gauge
\begin{equation}
\nabla\cdot\mathbf A+\frac1{c^2}\partial_t\phi=0
\quad\Longleftrightarrow\quad
\partial_\mu A^\mu=0,\qquad A^\mu=(\phi/c,\mathbf A),
\label{eq:mb-lorenz}
\end{equation}
the four equations decouple into four wave equations,
\begin{equation}
\Box A^\mu=\mu_0 J^\mu,\qquad J^\mu=(c\rho,\mathbf J).
\label{eq:mb-wave-A}
\end{equation}

Two remarks will matter later.

\begin{enumerate}
\item Applying $\partial_\mu$ to~\eqref{eq:mb-wave-A} and using
$\partial_\mu J^\mu=0$ gives $\Box(\partial_\mu A^\mu)=0$. The Lorenz
condition is therefore itself a \emph{homogeneous wave equation}: it holds
for all times if $\partial_\mu A^\mu$ and its time derivative vanish on one
initial surface. It is a condition on the data.
\item The Lorenz condition is preserved by a further gauge transformation
if and only if $\Box\chi=0$. The residual gauge freedom is itself a
homogeneous solution of the wave equation.
\end{enumerate}

So even the ambiguity in the description of the field lives in the space of
source-free solutions. This is the first hint that the homogeneous equation
is not an afterthought.

%----------------------------------------------------------------------
\section{Anatomy of the solutions: particular and homogeneous}
\label{sec:mb-anatomy}

\subsection{Superposition and the Laplace analogy}

Because~\eqref{eq:mb-wave-A} is linear, every solution is
\begin{equation}
A^\mu=A^\mu_{\rm p}+A^\mu_{\rm h},\qquad
\Box A^\mu_{\rm p}=\mu_0J^\mu,\qquad \Box A^\mu_{\rm h}=0 .
\label{eq:mb-split}
\end{equation}
The particular part is generated by the sources. The homogeneous part is
generated by nothing in the region under consideration.

An analogy from electrostatics is instructive. A solution of
$\nabla^2\phi=0$ in a region $\Omega$ vanishes identically if $\phi=0$ on
$\partial\Omega$ (maximum principle); the uniform field inside a capacitor
exists only because the plates carry non-zero boundary values. A free
solution is the shadow of its boundary data. The wave equation obeys the
same logic, with one difference: the ``boundary'' now includes the initial
surface, and since the equation is second order in time the data on it
consist of two functions (the field and its time derivative; for Maxwell's
equations, $\mathbf E$ and $\mathbf B$).

\subsection{Green functions}

With the convention $\Box G=\delta(t)\,\delta^3(\mathbf x)$ there are two
natural Green functions,
\begin{equation}
G_{\rm ret}(\mathbf x,t)=\frac{\delta(t-r/c)}{4\pi r},\qquad
G_{\rm adv}(\mathbf x,t)=\frac{\delta(t+r/c)}{4\pi r},
\qquad r=|\mathbf x| .
\label{eq:mb-green}
\end{equation}
Each gives a particular solution, $A^\mu_{\rm ret/adv}=\mu_0\,G_{\rm ret/adv}*J^\mu$.
Their difference
\begin{equation}
G_{\rm ret}-G_{\rm adv}=\frac{\delta(t-r/c)-\delta(t+r/c)}{4\pi r}
\label{eq:mb-pauli-jordan}
\end{equation}
satisfies $\Box(G_{\rm ret}-G_{\rm adv})=0$ everywhere, including at $r=0$,
where it is regular as a distribution (it behaves as $-\delta'(t)/2\pi c$).
Hence $A_{\rm ret}-A_{\rm adv}$ is a perfectly good homogeneous solution,
generated by the sources but source-free as a field. We will meet it again
in Section~\ref{sec:mb-arrow}.

\subsection{Plane waves}

In a region free of charges one may choose $\phi=0$, $\nabla\cdot\mathbf A=0$,
and the general homogeneous solution is
\begin{equation}
\begin{split}
\mathbf A_{\rm h}(\mathbf x,t)=\int\!\frac{d^3k}{(2\pi)^3}\sum_{\lambda=1,2}
\Big[a_\lambda(\mathbf k)\,&\boldsymbol\varepsilon_\lambda(\mathbf k)\,
e^{i(\mathbf k\cdot\mathbf x-\omega_k t)}+\text{c.c.}\Big],\\
&\omega_k=c|\mathbf k|,\qquad
\mathbf k\cdot\boldsymbol\varepsilon_\lambda=0 .
\end{split}
\label{eq:mb-modes}
\end{equation}
There are two polarizations per wave vector. This is the content of the
word ``photon'' at the classical level: the amplitudes $a_\lambda(\mathbf k)$
are free data, unconstrained by any source.

\subsection{No wave without data}

\begin{quote}
\textbf{Proposition 1 (no wave without data).} Let $\mathbf E,\mathbf B$
solve the source-free Maxwell equations near the surface $t=t_0$. If
$\mathbf E=\mathbf B=0$ on the ball $|\mathbf x-\mathbf x_0|<R$ at $t=t_0$,
then $\mathbf E=\mathbf B=0$ throughout the cone
$|\mathbf x-\mathbf x_0|<R-c|t-t_0|$.
\end{quote}

\emph{Proof.} Let $u=\tfrac12(\varepsilon_0E^2+B^2/\mu_0)$ and
$\mathbf S=\mathbf E\times\mathbf B/\mu_0$. Poynting's theorem for a ball
whose radius shrinks at speed $c$ gives
\begin{equation}
\frac{d}{dt}\int_{|\mathbf x-\mathbf x_0|<R-c(t-t_0)}\!\! u\,d^3x
=-\oint\big(\mathbf S\cdot\mathbf n+c\,u\big)\,d\sigma\le0,
\end{equation}
since $|\mathbf S|\le c\,u$. The energy is non-negative and starts at zero,
so it stays zero. $\square$

The consequence is exactly the statement of the introduction. A non-trivial
homogeneous solution is one whose Cauchy data $(\mathbf E,\mathbf B)$ do
not vanish somewhere on every spacelike slice; equivalently, in a
scattering situation, one that carries incoming radiation from past
infinity, or outgoing radiation to future infinity. In particular, ``no
incoming radiation'' is a \emph{boundary condition}, and it is the one
that turns the retarded solution into the solution of a physical problem.

\medskip
\noindent
\textbf{Caveat on what counts as data.} Proposition~1 concerns Cauchy data
(the field \emph{and} its time derivative on one slice). Data of a different
kind, for example the field alone on two different time slices, do not
suffice; we will see in Section~\ref{sec:mb-example} a non-trivial
solution that vanishes on the \emph{entire} boundary of a space-time box.
This distinction turns out to be central.

%----------------------------------------------------------------------
\section{The action principle and the term we discard}
\label{sec:mb-action}

From here on we use Heaviside--Lorentz units with $c=\varepsilon_0=\mu_0=1$
and the conventions listed at the head of the chapter.

\subsection{The variation of the action}

Take
\begin{equation}
S[A]=\int_V d^4x\Big[-\tfrac14F_{\mu\nu}F^{\mu\nu}-J^\mu A_\mu\Big],
\qquad F_{\mu\nu}=\partial_\mu A_\nu-\partial_\nu A_\mu ,
\label{eq:mb-action}
\end{equation}
on a space-time region $V$. Since $\delta(-\tfrac14F^2)=-F^{\mu\nu}\partial_\mu\delta A_\nu$,
integration by parts gives
\begin{equation}
\delta S=\int_V d^4x\,\big(\partial_\mu F^{\mu\nu}-J^\nu\big)\,\delta A_\nu
\;-\;\oint_{\partial V}d\Sigma_\mu\,F^{\mu\nu}\,\delta A_\nu .
\label{eq:mb-variation}
\end{equation}
The first term yields the field equations $\partial_\mu F^{\mu\nu}=J^\nu$.
The second is the total divergence we ``neglect''. For
$V=[t_1,t_2]\times\Omega$ it splits into a contribution from the two time
slices and one from the spatial boundary:
\begin{equation}
\oint_{\partial V}\!d\Sigma_\mu F^{\mu\nu}\delta A_\nu
=\Big[\int_\Omega d^3x\;\mathbf E\cdot\delta\mathbf A\Big]_{t_1}^{t_2}
+\int_{t_1}^{t_2}\!dt\oint_{\partial\Omega}\!d\sigma\;
\mathbf n\cdot\big(\mathbf E\,\delta\phi+\delta\mathbf A\times\mathbf B\big).
\label{eq:mb-boundary-split}
\end{equation}
In the language of Hamiltonian mechanics, the first piece is
$-\int\boldsymbol\pi\cdot\delta\mathbf A$ with $\boldsymbol\pi=-\mathbf E$ the
momentum conjugate to $\mathbf A$: the boundary term is ``momentum times
variation of the coordinate'', as in $[p\,\delta q]_{t_1}^{t_2}$ in
particle mechanics. In the covariant multiple-integral (De~Donder--Weyl)
language it is the pairing $\pi^{\mu\nu}\delta A_\nu$ with the covariant
momenta $\pi^{\mu\nu}=\partial L/\partial(\partial_\mu A_\nu)=-F^{\mu\nu}$.
The field strength \emph{is} the momentum.

\subsection{What ``$\delta A=0$ on the boundary'' actually assumes}

The usual argument sets $\delta A_\nu=0$ on $\partial V$ and the surface
term disappears. Three observations on this step.

\begin{enumerate}
\item \textbf{What vanishes is the variation, not the field.} The
background field $A$ may be non-zero on $\partial V$. The term
$\mathbf E\cdot\delta\mathbf A$ vanishes because $\delta\mathbf A=0$, not
because $\mathbf E=0$. The photon content, if any, sits in $\mathbf E$ and
is invisible to this step.

\item \textbf{Fixing $A$ at both ends is a two-point problem.} Physical
problems are initial-value problems: we know the past and want the future.
Holding $A$ fixed at $t_1$ \emph{and} $t_2$ asks for something we do not
know in advance, and the two-point problem for a wave equation need not
have a solution, nor a unique one (Section~\ref{sec:mb-example}). In
particle mechanics the same issue is hidden by the fact that Hamilton's
principle is used only to \emph{derive} the equations of motion.

\item \textbf{Spatial infinity is subtler than it looks.} For sources that
were switched on at a finite time the retarded field vanishes at fixed $t$
as $r\to\infty$ (the radiation has not arrived yet), so the spatial term
vanishes on every finite time slab. But radiation falls off only as $1/r$
while the area of a sphere grows as $r^2$, so surface integrals of
quadratic expressions such as $\delta\mathbf A\times\mathbf B$ over a
sphere in the radiation zone are $O(1)$, not zero. The energy that leaves
the system crosses precisely such spheres.
\end{enumerate}

\subsection{On shell, the free action is a pure divergence}

The identity $F_{\mu\nu}F^{\mu\nu}=2F^{\mu\nu}\partial_\mu A_\nu$ gives
\begin{equation}
-\tfrac14F^2=-\tfrac12\partial_\mu\big(F^{\mu\nu}A_\nu\big)
+\tfrac12\big(\partial_\mu F^{\mu\nu}\big)A_\nu .
\end{equation}
On a solution, $\partial_\mu F^{\mu\nu}=J^\nu$, so
\begin{equation}
S_{\rm on\text{-}shell}
=-\tfrac12\int_V d^4x\,J^\mu A_\mu
-\tfrac12\oint_{\partial V}d\Sigma_\mu\,F^{\mu\nu}A_\nu .
\label{eq:mb-onshell}
\end{equation}
For the free field ($J=0$) the on-shell action is \emph{entirely} a
boundary term. If we discard divergences in the action, every free solution
has the same action, namely zero. The first term in~\eqref{eq:mb-onshell}
is the familiar ``half of $-J\!\cdot\!A$'' of the Coulomb interaction energy;
with $A=A_{\rm p}+A_{\rm h}$ it contains $A_{\rm h}$ only through its
coupling to the sources (exchange of energy with free radiation). The
\emph{self}-action of the free radiation is a pure divergence.

Note the two different divergences in play, and the two different
neglects they correspond to:
\begin{itemize}
\item the surface term in $\delta S$, eq.~\eqref{eq:mb-variation}, which
we drop to derive the field equations (harmless for the equations, but it
discards the symplectic potential, see Section~\ref{sec:mb-symplectic});
\item the surface term in $S$ itself, eq.~\eqref{eq:mb-onshell}, which is
dropped whenever one evaluates the action on solutions (this discards the
free-field part completely).
\end{itemize}

\subsection{Worked example: a standing wave in a cavity}
\label{sec:mb-example}

Take a single polarization between perfectly conducting walls at $x=0$ and
$x=\ell$. Then $A\equiv A_y(x,t)$, the Lagrangian per unit transverse area
is $L=\tfrac12(A_t^2-A_x^2)$, and the boundary condition is $A=0$ at the
walls. Write $\Box=\partial_t^2-\partial_x^2$. Then
\begin{equation}
L=\tfrac12\partial_\mu\big(A\,\partial^\mu A\big)-\tfrac12A\,\Box A
=\tfrac12\big[\partial_t(AA_t)-\partial_x(AA_x)\big]\quad\text{on shell.}
\end{equation}
Consider the mode $A=a\sin(kx)\cos(\omega t+\delta)$ with
$\omega=k=n\pi/\ell$. The $x$-divergence integrates to zero (walls), and
\begin{equation}
S=\int_0^T\!dt\int_0^\ell\!dx\,L
=\tfrac12\Big[\int_0^\ell\!dx\,AA_t\Big]_0^T
=-\frac{a^2\omega\ell}{8}\Big[\sin(2\omega T+2\delta)-\sin2\delta\Big].
\label{eq:mb-cavity-action}
\end{equation}
One checks by direct integration of $\tfrac12(A_t^2-A_x^2)$ that this is
correct. Three lessons follow.

\begin{enumerate}
\item The on-shell action of a perfectly healthy, non-trivial free
solution is non-zero and consists of a boundary term at $t=0$ and $t=T$
only. Dropping divergences would assign it the value $0$.
\item Taking $\delta=\pi/2$ gives $A=-a\sin(kx)\sin(\omega t)$, which
vanishes at $t=0$. It also vanishes at $t=T$ whenever $\omega T=m\pi$. For
$T=m\ell/n$ (any rational multiple of $\ell$) there is a non-trivial
solution with $A=0$ on the \emph{whole boundary} of the space-time
rectangle. The fixed-endpoint variational problem is degenerate there: the
second variation has a null direction. This is the field-theoretic analogue
of a conjugate point in Jacobi's theory.
\item The resolution of the apparent paradox is that fixing $A$ is only
half of the Cauchy data. Here $A_t=-\omega a\sin kx\cos\omega t\neq0$ at
$t=0$: the momentum $\mathbf E$ is non-zero, and it multiplies
$\delta A=0$ in the discarded term. The photon is hiding in the factor that
multiplies the vanishing variation.
\end{enumerate}

\subsection[Wave equation from a shifted Lagrangian]{The same wave equation from a Lagrangian that differs by a divergence}

The identity
\begin{equation}
\begin{split}
-\tfrac14F_{\mu\nu}F^{\mu\nu}
&=-\tfrac12\partial_\mu A_\nu\partial^\mu A^\nu+\tfrac12(\partial_\mu A^\mu)^2
+\partial_\mu K^\mu,\\
K^\mu&=\tfrac12\big(A_\nu\partial^\nu A^\mu-A^\mu\partial_\nu A^\nu\big),
\end{split}
\label{eq:mb-feynman}
\end{equation}
shows that Maxwell's Lagrangian plus the gauge-fixing term
$-\tfrac12(\partial_\mu A^\mu)^2$ equals $-\tfrac12\partial_\mu
A_\nu\partial^\mu A^\nu$ up to the divergence $\partial_\mu K^\mu$. After
discarding $\partial_\mu K^\mu$, the Lagrangian
\begin{equation}
L_{\rm F}=-\tfrac12\partial_\mu A_\nu\,\partial^\mu A^\nu-J^\mu A_\mu
\end{equation}
gives $\Box A^\nu=J^\nu$ directly, four decoupled wave equations. This is
the cleanest derivation of the wave equation from an action, and it is made
possible by dropping a divergence. The price is that the Lorenz condition
is no longer implied by the equations of motion; it must be imposed on the
free solutions, where it is a homogeneous wave equation with its own
boundary data (Section~\ref{sec:mb-wave}, remark~1), exactly the
Gupta--Bleuler condition of the quantum theory.

%----------------------------------------------------------------------
\section[The discarded term: symplectic structure]{What the discarded term contains: symplectic structure and photons}
\label{sec:mb-symplectic}

The boundary term at a time slice is
$\Theta[\delta A]=-\int d^3x\,\mathbf E\cdot\delta\mathbf A$. Its exterior
derivative on the space of solutions is the symplectic form. Concretely,
for two source-free solutions $A_1,A_2$ define the current
\begin{equation}
j^\mu[A_1,A_2]=F_1^{\mu\nu}A_{2\nu}-F_2^{\mu\nu}A_{1\nu}.
\end{equation}
Using $\partial_\mu F^{\mu\nu}=0$ and $F^{\mu\nu}\partial_\mu A_\nu=\tfrac12F^{\mu\nu}F_{\mu\nu}$,
\begin{equation}
\partial_\mu j^\mu=\tfrac12F_1^{\mu\nu}F_{2\mu\nu}-\tfrac12F_2^{\mu\nu}F_{1\mu\nu}=0 .
\end{equation}
The charge of this current,
\begin{equation}
\Omega(A_1,A_2)=\int d^3x\,\big(\mathbf E_1\cdot\mathbf A_2-\mathbf E_2\cdot\mathbf A_1\big),
\label{eq:mb-symplectic}
\end{equation}
is therefore independent of the slice. It is gauge invariant for
$\nabla\cdot\mathbf E=0$, because $\Omega(A,\nabla\chi)=-\int\chi\,\nabla\cdot\mathbf E=0$.
Up to a convention-dependent sign, $\Omega$ is the exterior derivative of
the boundary term. The boundary term we discard in~\eqref{eq:mb-variation}
\emph{is} the symplectic potential, and the conserved symplectic form is
what remains of it after one more variation. This is the structure on which
canonical quantization rests (the ``covariant phase space'' of Crnkovi\'c--Witten
and Lee--Wald).

Evaluate $\Omega$ on the plane waves $A_{\mathbf k,\lambda}=\boldsymbol\varepsilon_\lambda e^{i(\mathbf k\cdot\mathbf x-\omega t)}$.
A short computation gives
\begin{equation}
\Omega\big(A_{\mathbf k\lambda},A^{*}_{\mathbf k'\lambda'}\big)
=2i\omega_k\,(2\pi)^3\delta^3(\mathbf k-\mathbf k')\,\delta_{\lambda\lambda'},
\qquad
\Omega\big(A_{\mathbf k\lambda},A_{\mathbf k'\lambda'}\big)=0 .
\end{equation}
In the mode amplitudes $a_\lambda(\mathbf k)$ of~\eqref{eq:mb-modes} the
symplectic form is therefore diagonal, so $a_\lambda$ and $a_\lambda^*$ are
canonically conjugate up to normalization. Quantization replaces them by
annihilation and creation operators. For one mode in a box of volume $V$
the time-averaged field energy is $2\omega^2|a|^2V$, which gives
\begin{equation}
N=\frac{2\,\omega\,V\,|a|^2}{\hbar}
\end{equation}
quanta of energy $\hbar\omega$. The photon number is therefore read off
from $\Omega$, that is, from the divergence term of the action, and from
nowhere else.

It is worth stressing that this is more than the on-shell value of $S$. For
a traveling plane wave in a periodic box, $E^2=B^2$ pointwise, so
$L=\tfrac12(E^2-B^2)=0$ and the on-shell action vanishes even with the
boundary term kept (Exercise~\ref{ex:mb-plane}). The wave nevertheless
carries photons. The information is not in the \emph{value} of the
boundary term but in its \emph{variation}.

%----------------------------------------------------------------------
\section{The action cannot tell past from future}
\label{sec:mb-arrow}

Insert the solution $A^\mu=G*J^\mu$ of $\Box A^\mu=J^\mu$ into the first
term of~\eqref{eq:mb-onshell}:
\begin{equation}
S_{\rm eff}[J]=-\tfrac12\int\!\!\int d^4x\,d^4y\;J^\mu(x)\,G(x-y)\,J_\mu(y).
\end{equation}
The integrand is symmetric under $x\leftrightarrow y$, so only the
symmetric part of $G$ contributes. The retarded and advanced Green
functions give the \emph{same} $S_{\rm eff}$, equal to the one built from
\begin{equation}
G_{\rm sym}=\tfrac12\big(G_{\rm ret}+G_{\rm adv}\big).
\end{equation}
This is the origin of the Fokker--Schwarzschild--Tetrode action: a
variational principle for charges that interact directly, time-symmetrically.
Nothing in it distinguishes retarded from advanced.

The difference between the retarded solution and the time-symmetric one is
\begin{equation}
A_{\rm ret}=A_{\rm sym}+A_{\rm rad},\qquad
A_{\rm rad}=\tfrac12\big(A_{\rm ret}-A_{\rm adv}\big),
\label{eq:mb-dirac}
\end{equation}
and by~\eqref{eq:mb-pauli-jordan} $A_{\rm rad}$ is a source-free solution,
regular on the worldline of the source. It is non-zero, and it is
homogeneous: exactly the kind of object that the action sees only through
boundary terms. Dirac showed that evaluated on the worldline of a point
charge it yields the radiation-reaction force, which for slow motion is
\begin{equation}
\mathbf F_{\rm rr}=\frac{q^2}{6\pi}\,\frac{d^3\mathbf x}{dt^3}
\qquad(\text{Heaviside--Lorentz, }c=1),
\end{equation}
the Abraham--Lorentz force. The arrow of time in electrodynamics thus does
not come from the action. It comes from the choice $A_{\rm h}=A_{\rm in}=0$
in the past, which is a boundary condition on the homogeneous solution. The
Wheeler--Feynman absorber theory attempts to derive that condition from the
requirement of complete absorption in the future, again a statement about
boundary data.

\subsection*{A remark on the two worlds}

In the terminology of this book the situation is a clean example of two
worlds that should not be mixed. Causality, ``what has actually been
radiated'' and ``what we see in a detector'' belong to the real world. The
variational principle, with its fixed endpoints and discarded divergences,
is a construction of the ideal world: time-symmetric, and blind to
boundary data by design. If one asks the ideal-world construction to
reproduce a real-world fact (the retarded, causal behavior of radiation, or
the number of photons in a beam), it cannot. The information was removed in
the step ``neglect the divergence'', and has to be put back in by hand as a
boundary condition. The pattern is the same as in the other examples of this book: a
rigorous derivation inside the ideal world is asked a question that belongs
to the real world.

The remedy is not to abandon the action but to keep track of the term
that was discarded. Modern formulations do this in different ways. The
covariant phase space keeps the symplectic potential. The closed-time-path
(Schwinger--Keldysh) formulation, applied to classical mechanics by Galley,
doubles the variables so that the action can be varied with initial data
only, and the boundary term at the final time cancels between the two
copies. Both are ways of declining to discard the divergence.

%----------------------------------------------------------------------
\section[Other discarded divergences]{Other places where discarded divergences matter}
\label{sec:mb-others}

\begin{itemize}
\item \textbf{Stress tensor.} The canonical energy--momentum tensor of the
Maxwell field is not symmetric, and the Belinfante symmetric tensor differs
from it by a divergence. The total four-momentum agrees if the surface term
vanishes. For radiation, which falls only as $1/r$, the corresponding
surface integrals for \emph{angular momentum} do not vanish automatically,
and the split into ``spin'' and ``orbital'' parts of the angular
momentum of light is sensitive to which surface terms are discarded.
\item \textbf{The $\theta$-term.} The Lagrangian $\theta\,F_{\mu\nu}\tilde F^{\mu\nu}$ is a pure divergence,
\begin{equation*}
F_{\mu\nu}\tilde F^{\mu\nu}=2\,\partial_\mu\big(\varepsilon^{\mu\nu\rho\sigma}A_\nu\partial_\rho A_\sigma\big),
\end{equation*}
proportional to $\mathbf E\cdot\mathbf B$ (the sign depends on the convention
for $\varepsilon^{0123}$). It changes no local field equation. It matters,
however, for non-trivial boundary data or topology (Chern--Simons terms,
instantons), where the divergence cannot be integrated away.
\item \textbf{Gravity.} The Einstein--Hilbert action is second order in
derivatives of the metric; the boundary term that makes the variational
principle well posed with fixed boundary metric (Gibbons--Hawking--York) is
a divergence that cannot be dropped.
\end{itemize}

%----------------------------------------------------------------------
\section{Summary}
\label{sec:mb-summary}

\begin{enumerate}
\item The wave equation is a consequence of Maxwell's equations. Its
solutions split as $A=A_{\rm p}+A_{\rm h}$ into a part driven by the
sources and a source-free part.
\item By Proposition~1, a non-trivial homogeneous solution is
non-trivial because its Cauchy data, or its data at infinity, are
non-zero. Free radiation is the boundary data of the problem.
\item In the action principle the boundary data enter only through
divergences: the boundary term of $\delta S$ (the symplectic potential,
$\mathbf E\cdot\delta\mathbf A$) and the boundary term of the on-shell
action (the whole of the free-field action).
\item ``Neglecting the divergence'' is harmless for deriving the field
equations, but it is equivalent to switching off the physics carried by the
homogeneous solutions: the symplectic structure and hence the photon
number, the distinction between retarded and advanced solutions, and
radiation reaction.
\item The statement ``$\delta A=0$ on the boundary'' fixes only half of
the Cauchy data. The other half, $\mathbf E$, multiplies the vanishing
variation, and this is where the photons are hiding.
\end{enumerate}

%----------------------------------------------------------------------
\section*{Exercises}

\begin{enumerate}
\item Show that the Lorenz condition, once satisfied together with its
time derivative on an initial surface, holds for all times when
$\partial_\mu J^\mu=0$. Show that the residual gauge transformations
satisfy $\Box\chi=0$.
\item Verify that $\partial_\mu j^\mu=0$ for the symplectic current
and that $\Omega(A,\nabla\chi)=0$ when $\nabla\cdot\mathbf E=0$.
\item Verify~\eqref{eq:mb-cavity-action} by integrating $\tfrac12(A_t^2-A_x^2)$
directly. For which $T$ does the Dirichlet problem $A=0$ on the whole
boundary of $[0,\ell]\times[0,T]$ have a non-trivial solution? Relate your
answer to Jacobi's conjugate points.
\item\label{ex:mb-plane} Show that for a circularly polarized plane wave
in a periodic box the Lagrangian density vanishes identically, hence so
does the on-shell action including the boundary term, while the symplectic norm
$i\,\Omega(A^{(-)},A^{(+)})$ of its positive-frequency part $A^{(+)}$ is non-zero.
\item Derive the identity~\eqref{eq:mb-feynman} and show that the Lorenz
condition $\partial_\mu A^\mu=0$ follows from $\Box A^\mu=J^\mu$ only up to
a homogeneous solution.
\item Show that $G_{\rm ret}-G_{\rm adv}$ is regular at $r=0$, and
compute the force $\mathbf F_{\rm rr}$ on a slowly moving charge from
$A_{\rm rad}=\tfrac12(A_{\rm ret}-A_{\rm adv})$.
\end{enumerate}

%----------------------------------------------------------------------
\section*{Notes and further reading}

The treatment of the wave equation and the Green functions follows the
standard accounts in Jackson and in Landau and Lifshitz (Vol.~2). The
decomposition~\eqref{eq:mb-dirac} and the radiation-reaction force are due
to Dirac (1938); the time-symmetric action and the absorber condition to
Wheeler and Feynman (1945, 1949). The covariant phase space is described by
Crnkovi\'c and Witten (1987) and Lee and Wald (1990). Galley (2013) gives
the initial-value formulation of the action principle with doubled
variables.
% TODO: add \ref to the De Donder--Weyl / multiple-integral chapter
% (Sec. 4.1: covariant momenta pi^{mu nu} = -F^{mu nu}).
